\documentclass[12pt,a4paper,oneside]{ctexart}
\usepackage{amsmath,amsthm,amssymb,bm,graphicx,hyperref,mathrsfs}

\title{\textbf{作业}}
% \author{黄伟宁}
\date{\today}
\linespread{1.5}
\newcounter{problemname}
\newenvironment{problem}{\stepcounter{problemname}\par\noindent\textbf{题目\arabic{problemname}. }}{\\\par}
\newenvironment{solution}{\par\noindent\textbf{解答. }}{\\\par}
\newenvironment{note}{\par\noindent\textbf{题目\arabic{problemname}的注记. }}{\\\par}
\usepackage{geometry}
\geometry{left=2.54cm,right=2.54cm,top=2.9cm,bottom=2.9cm}

\begin{document}

\maketitle
% \begin{problem}
% 这里是题目.
% \end{problem}

% \begin{solution}
%     这里是解答.
% \end{solution}

% \begin{note}
%     这里是注记.
% \end{note}

% \begin{problem}
% \end{problem}
\begin{problem}
    设$ a>b>0 $,求$ a+\frac{1}{b(a-b)} $的最小值.
\end{problem}

\begin{problem}
    设$ a,b,c $为正实数,若$ M=2(\frac{a+b}{2}-\sqrt{ab}),N=3(\frac{a+b+c}{3}-\sqrt[3]{abc}) $,试比较$ M,N $的大小.
\end{problem}

\begin{problem}
    设$ a,b,c $为正实数,证明:

    (1)$ (a+b+c)(\frac{1}{a}+\frac{1}{b}+\frac{1}{c})\geqslant 9 $;

    (2)$ (a+b+c)(\frac{1}{a+b}+\frac{1}{b+c}+\frac{1}{c+a})\geqslant \frac{9}{2} $;

    (3)$ \frac{a}{b+c}+\frac{b}{c+a}+\frac{c}{a+b}\geqslant \frac{3}{2} $.
\end{problem}

\begin{problem}
    设$ a,b,c $为正数,证明:$ \frac{a^2}{b}+\frac{b^2}{c}+\frac{c^2}{a}+2(a+b+c)\geqslant \frac{(a+b+c)^3}{ab+bc+ca} $.
\end{problem}

\begin{problem}
    设$ a,b,c,d>0,a+b+c+d=1 $,证明:$ (1-\sqrt{a})(1-\sqrt{b})(1-\sqrt{c})(1-\sqrt{d})\geqslant \sqrt{abcd} $.
\end{problem}

\begin{problem}
    若$ a,b,c $为非负实数,且$ a^2+b^2+c^2-ab-bc-ca=25 $,则$ a+b+c $的最小值为\underline{\qquad\qquad}.
\end{problem}

\begin{problem}
    若$ a>b>c $且$ 3a+2b+c=0 $,则$ \frac{b+c}{a} $的取值范围是\underline{\qquad\qquad}.
\end{problem}

\begin{problem}
    设$ 1\leqslant a\leqslant b\leqslant c\leqslant d\leqslant 4 $,则$ \frac{a}{b}+\frac{c}{d} $的最小值为\underline{\qquad\qquad}.
\end{problem}

\begin{problem}
    已知$ \triangle ABC $的三边长$ a,b,c $满足$ b+c\leqslant 2a,c+a\leqslant 2b $,则$ \frac{b}{a} $的取值范围是\underline{\qquad\qquad}.
\end{problem}

\end{document}